\documentclass[10pt,a4paper]{amsart}
\usepackage[margin=19mm]{geometry}
\usepackage{amsmath,amssymb,mathtools}
\usepackage[scr=rsfs,cal=euler]{mathalfa}
\usepackage{xfrac}
\usepackage[unicode,hidelinks,bookmarks=false]{hyperref}
\newcommand{\ie}{i.e.\ }
\newcommand{\cf}{cf.\ }
\newcommand{\resp}{respectively\ }
\newcommand{\Subsection}{\subsection}
\providecommand{\nobreakdash}{\nobreak\hbox{-}}
\setlength{\emergencystretch}{2em}
\multlinegap0pt
\usepackage{mathtools}
\usepackage{mathrsfs}
\usepackage{latexsym}
\usepackage{textcomp}
\usepackage{array}
\usepackage{booktabs}
\usepackage{enumitem}
\usepackage{setspace}
\usepackage{tikz}
\usepackage{tikz-cd}
\usetikzlibrary{trees}
\usepackage{amscd}
\usepackage{times}
\hypersetup{hidelinks}
\setlength{\textheight}{210mm}
\setlength{\topmargin}{0.46cm}
\setlength{\textwidth}{152mm}
\setlength{\evensidemargin}{0.60cm}
\setlength{\oddsidemargin}{0.60cm}
\setcounter{MaxMatrixCols}{30}
\numberwithin{equation}{section}
\setcounter{tocdepth}{1}
\renewcommand{\sectionmark}[1]{\markright{\thesection \ #1}{}}
\newtheorem{Theorem}{Theorem}[section]
\newtheorem{theoremn}{Theorem}
\newtheorem{Lemma}[Theorem]{Lemma}
\newtheorem{Proposition}[Theorem]{Proposition}
\newtheorem{Corollary}[Theorem]{Corollary}
\newtheorem{Conjecture}[Theorem]{Conjecture}
\newtheorem*{question}{Question}
\newtheorem*{remark}{Remark}
\newtheorem*{conjecturen}{Conjecture}
\newtheorem{Definition}[Theorem]{Definition}
\newtheorem{Construction}[Theorem]{Construction}
\newtheorem{Notation}[Theorem]{Notation}
\newtheorem{Example}[Theorem]{Example}
\newtheorem{Remark}[Theorem]{Remark}
\newtheorem{Question}[Theorem]{Question}
\newtheorem{Assumption}[Theorem]{Assumption}
\newtheorem{say}[Theorem]{}
\newcommand{\arXiv}[1]{\href{https://arxiv.org/abs/#1}{arXiv:#1}}
\newcommand{\ndo}{\node[draw,circle,inner sep=1]}
\newcommand{\Spec}{\operatorname{Spec}}
\newcommand{\Hom}{\operatorname{Hom}}
\newcommand{\Spf}{\operatorname{Spf}}
\newcommand{\Proj}{\operatorname{Proj}}
\newcommand{\Sing}{\operatorname{Sing}}
\newcommand{\Supp}{\operatorname{Supp}}
\newcommand{\Exc}{\operatorname{Exc}}
\newcommand{\codim}{\operatorname{codim}}
\newcommand{\ord}{\operatorname{ord}}
\newcommand{\val}{\operatorname{val}}
\newcommand{\mld}{\operatorname{mld}}
\newcommand{\tmldtor}{\operatorname{tmld}_{\mathrm{tor}}}
\newcommand{\lct}{\operatorname{lct}}
\newcommand{\tlct}{\operatorname{tlct}}
\newcommand{\lcodim}{\operatorname{lcodim}}
\newcommand{\Diff}{\operatorname{Diff}}
\newcommand{\Cont}{\operatorname{Cont}}
\newcommand{\Nklt}{\operatorname{Nklt}}
\newcommand{\Nlc}{\operatorname{Nlc}}
\newcommand{\TNklt}{\operatorname{Nklt}_{\tan}}
\newcommand{\Jac}{\operatorname{Jac}}
\newcommand{\Nash}{\operatorname{Nash}}
\newcommand{\coeff}{\operatorname{coeff}}
\newcommand{\Cone}{\operatorname{Cone}}
\newcommand{\Res}{\operatorname{Res}}
\newcommand{\Aut}{\operatorname{Aut}}
\newcommand{\Ext}{\operatorname{Ext}}
\newcommand{\Hilb}{\operatorname{Hilb}}
\newcommand{\Ker}{\operatorname{Ker}}
\newcommand{\Isom}{\operatorname{Isom}}
\newcommand{\im}{\operatorname{Im}}
\newcommand{\Sym}{\operatorname{Sym}}
\newcommand{\Pic}{\operatorname{Pic}}
\newcommand{\Bs}{\operatorname{Bs}}
\newcommand{\Bl}{\operatorname{Bl}}
\newcommand{\Gr}{\operatorname{Gr}}
\newcommand{\End}{\operatorname{End}}
\newcommand{\MJ}{\mathrm{MJ}}
\newcommand{\sn}{\mathrm{sn}}
\newcommand{\tanMJ}{\mathrm{tan\text{-}MJ}}
\newcommand{\Sigmatan}{\Sigma_{\mathrm{tan}}}
\newcommand{\bbound}{\mathfrak b}
\newcommand{\Adm}{\mathsf{Adm}}
\newcommand{\SN}{\mathsf{SN}}
\newcommand{\red}{\mathrm{red}}
\newcommand{\Strata}{\mathsf{Str}}
\newcommand{\Cyl}{\mathsf{Cyl}}
\newcommand{\inv}{\mathrm{inv}}
\newcommand{\tmld}{\operatorname{tmld}}
\newcommand{\Jetm}{J_m}
\newcommand{\dd}{\mathrm d}
\newcommand{\eps}{\varepsilon}
\newcommand{\lra}{\longrightarrow}
\newcommand{\Rlambda}{\mathcal R_\lambda}
\newcommand{\AF}{A_{\mathcal F}}
\newcommand{\Dinv}{D_{\mathrm{inv}}}
\newcommand{\bA}{\mathbb A}
\newcommand{\bC}{\mathbb C}
\newcommand{\bN}{\mathbb N}
\newcommand{\bQ}{\mathbb Q}
\newcommand{\bR}{\mathbb R}
\newcommand{\bZ}{\mathbb Z}
\newcommand{\PP}{\mathbb P}
\newcommand{\CC}{\mathbb C}
\newcommand{\ZZ}{\mathbb Z}
\newcommand{\cF}{\mathcal F}
\newcommand{\cG}{\mathcal G}
\newcommand{\cO}{\mathcal O}
\newcommand{\cB}{\mathcal B}
\newcommand{\cC}{\mathcal C}
\newcommand{\cJ}{\mathcal J}
\newcommand{\cN}{\mathcal N}
\newcommand{\cS}{\mathcal S}
\newcommand{\cT}{\mathcal T}
\newcommand{\sA}{\mathscr A}
\newcommand{\sB}{\mathscr B}
\newcommand{\sC}{\mathscr C}
\newcommand{\sD}{\mathscr D}
\newcommand{\sE}{\mathscr E}
\newcommand{\sF}{\mathscr F}
\newcommand{\sG}{\mathscr G}
\newcommand{\sH}{\mathscr H}
\newcommand{\sI}{\mathscr I}
\newcommand{\sJ}{\mathscr J}
\newcommand{\sK}{\mathscr K}
\newcommand{\sL}{\mathscr L}
\newcommand{\sM}{\mathscr M}
\newcommand{\sN}{\mathscr N}
\newcommand{\sO}{\mathscr O}
\newcommand{\sP}{\mathscr P}
\newcommand{\sQ}{\mathscr Q}
\newcommand{\sR}{\mathscr R}
\newcommand{\sS}{\mathscr S}
\newcommand{\sT}{\mathscr T}
\newcommand{\sU}{\mathscr U}
\newcommand{\sV}{\mathscr V}
\newcommand{\sW}{\mathscr W}
\newcommand{\sX}{\mathscr X}
\newcommand{\sY}{\mathscr Y}
\newcommand{\sZ}{\mathscr Z}
\newcommand{\GG}{\mathfrak G}
\newcommand{\jaco}{\mathfrak j}
\newcommand{\OO}{\mathcal O}
\newcommand{\Poiss}{\operatorname{Poisson}}
\newcommand{\Fol}{\operatorname{Fol}}
\newcommand{\ad}{\operatorname{ad}}
\newcommand{\rk}{\operatorname{rk}}
\newcommand{\slthree}{\mathfrak{sl}_3}
\newcommand{\sltwo}{\mathfrak{sl}_2}
\newcommand{\Spanop}{\operatorname{Span}}
\newcommand\Span[1]{\langle#1\rangle}
\newcommand{\Rat}{\textup{RatCurves}}
\DeclareMathOperator{\rank}{rank}
\def\bibaut#1{{\sc #1}}
\begin{document}
\title[Foliated and Mather-Jacobian discrepancies via tangential ar]{Foliated and Mather-Jacobian discrepancies via tangential arcs}
\author{Maurício Corrêa}
\date{}
\maketitle
\noindent\emph{Source and licence.} Foliated and Mather-Jacobian discrepancies via tangential arcs. Version arXiv:2607.01809v2 (CC BY 4.0). This material is distributed under the Creative Commons Attribution 4.0 International licence, http://creativecommons.org/licenses/by/4.0/, as recorded in the arXiv OAI licence metadata for this exact version and shown on the abstract page https://arxiv.org/abs/2607.01809v2. Full rights notice: licenses/arxiv-2607.01809-CC-BY-4.0.txt.\par\smallskip
\noindent\textit{Editorial scope.} Counted unit: the original \texttt{Theorem} environment \texttt{thm:carter-representability} at source lines 1024--1057 together with its complete proof at lines 1060--1202; the delivered document keeps source lines 999--1203 so that every referenced label and every citation resolves inside it. Native editable source is the paper's own concatenated TeX; the AMS wrapper is editorial formatting only. No publisher class, upstream script or unrelated artwork is redistributed.
\section{Formal tangential representability}
\label{sec:formal}

Carter's jet schemes of a foliation \cite{Carter,CarterThesis} provide the
finite-jet language motivating the tangential sector.  The discrepancy
calculations in what follows, however, use only the reduced infinite tangential arc
functor.  Thus, when a co-rank one foliation is locally generated by an
integrable one-form $\omega$, the reduced functor is used
\[
J_\infty^{\tan}(W,\cG;Z)_{\red}
=
\{\gamma\in J_\infty(W)_{\red}:\gamma(0)\in Z,\ \gamma^*\omega=0\}.
\]
This condition is independent of the chosen local generator, since
multiplying $\omega$ by a unit does not change the vanishing of its pull-back.
This is the precise arc-level consequence of Carter's tangent-jet viewpoint
which is needed in the sequel.  The finite tangent jet schemes may contain
nilpotent and truncation dependent structure; no equality of finite-level
Carter jet schemes is used in any cylinder or codimension calculation in this
article.

The next theorem is the only reduced Carter-type result required in what follows.  It is
best regarded as an arc-confinement statement rather than as a finite-jet
representability theorem.


\begin{Theorem} 
\label{thm:carter-representability}
Let $p\in W$ have a logarithmic simple adapted chart and write
\[
R=\bC[[x_1,x_2,x_3]],\qquad f=x_1\cdots x_r,
\qquad \Dinv=(f=0).
\]
Let
\[
\omega=
u\,x_1\cdots x_r
\left(
\sum_{i=1}^r \lambda_i\frac{dx_i}{x_i}
+
\sum_{j=r+1}^3 h_j(x)dx_j
\right)
\]
be the local generator, where $u$ is a unit and the residues satisfy the
positive non-resonance condition.  Then, for every closed subset
$Z\subset \Dinv$, the reduced tangential arc functor centred on $Z$
coincides with the reduced arc functor of $\Dinv$ centred on $Z$.
Equivalently, every reduced tangential formal arc centred on $Z$ factors
through the adapted invariant divisor $\Dinv$, and every arc on
$\Dinv$ is tangential.  In particular, for
$\Sigmatan=\Sing\cG\cap |\Dinv|$,
\[
J_\infty^{\tan}(W,\cG;\Sigmatan)_{\red}
=
J_\infty(\Dinv;\Sigmatan)_{\red}.
\]
This is an equality of reduced arc functors only.  No assertion is made about
the non-reduced finite Carter jet schemes, nor about equality of finite-level
reduced $m$-jet schemes for all $m$.
\end{Theorem}

\begin{proof}
The assertion is local at \(p\).  Since multiplication by a unit does not
change the vanishing of the pull-back of a one-form, we may omit \(u\).
Let
\[
\gamma:\Spec K[[t]]\longrightarrow W,
\qquad
\gamma(t)=(x_1(t),x_2(t),x_3(t)),
\]
be a formal arc centred on \(\Dinv\).  We first prove that a tangential arc
centred on \(\Dinv\) is contained in \(\Dinv\).

Assume, by contradiction, that \(\gamma\) is not contained in \(\Dinv\).
Then none of \(x_1(t),\ldots,x_r(t)\) is the zero series.  Write
\[
a_i=\ord_t x_i(t),\qquad
x_i(t)=t^{a_i}u_i(t),
\qquad
u_i(0)=c_i\in K^*,
\]
and set
\[
P(t)=x_1(t)\cdots x_r(t),\qquad
A=\ord_t P(t)=\sum_{i=1}^r a_i .
\]
Since the closed point of \(\gamma\) lies on \(\Dinv\), at least one \(a_i\)
is positive.  Thus \(a=(a_1,\ldots,a_r)\) is a non-zero element of
\(\bZ_{\geq 0}^r\).
Pulling back the logarithmic generator gives
\[
\gamma^*\omega
=
P(t)
\left(
\sum_{i=1}^r\lambda_i\frac{x_i'(t)}{x_i(t)}
+
\sum_{j=r+1}^3 h_j(\gamma(t))x_j'(t)
\right)dt .
\]
We analyse the two summands separately.  Since
\(x_i(t)=t^{a_i}u_i(t)\) with \(u_i(0)\neq0\), one has
\[
\frac{x_i'(t)}{x_i(t)}
=
\frac{a_i}{t}+\frac{u_i'(t)}{u_i(t)}.
\]
Hence
\[
\sum_{i=1}^r\lambda_i\frac{x_i'(t)}{x_i(t)}
=
\left(\sum_{i=1}^r a_i\lambda_i\right)\frac{1}{t}
+
\sum_{i=1}^r\lambda_i\frac{u_i'(t)}{u_i(t)}.
\]
The second term on the right-hand side is regular in \(t\).  Moreover
\[
P(t)
=
t^A\prod_{i=1}^r u_i(t)
=
t^A
\left(
\prod_{i=1}^r c_i+O(t)
\right).
\]
Therefore the logarithmic part contributes
\[
P(t)
\sum_{i=1}^r\lambda_i\frac{x_i'(t)}{x_i(t)}
=
\left(\sum_{i=1}^r a_i\lambda_i\right)
\left(\prod_{i=1}^r c_i\right)t^{A-1}
+
O(t^A).
\]
On the other hand, the regular part satisfies
\[
h_j(\gamma(t))x_j'(t)\in K[[t]]
\]
for every \(j>r\).  Hence
\[
P(t)\sum_{j=r+1}^3 h_j(\gamma(t))x_j'(t)
\in t^A K[[t]].
\]
Thus the coefficient of \(t^{A-1}dt\) in \(\gamma^*\omega\) is exactly
\[
\left(\sum_{i=1}^r a_i\lambda_i\right)
\prod_{i=1}^r c_i .
\]
By positive non-resonance,
\[
\sum_{i=1}^r a_i\lambda_i\neq0,
\]
and since each \(c_i\neq0\), this coefficient is non-zero.  Consequently
\(\gamma^*\omega\neq0\), contradicting tangency.  Hence every reduced
tangential arc centred on \(\Dinv\) is contained in \(\Dinv\).

Conversely, assume that \(\gamma\) factors through \(\Dinv\).  Since
$
\Dinv=(x_1\cdots x_r=0)
$
and \(K[[t]]\) is an integral domain, there exists an index
\(i_0\in\{1,\ldots,r\}\) such that
$
x_{i_0}(t)\equiv0.
$
Write \(f=x_1\cdots x_r\).  The logarithmic form is the regular one-form
\[
\omega
=
u\left(
\sum_{i=1}^r\lambda_i\frac{f}{x_i}\,dx_i
+
f\sum_{j=r+1}^3 h_j\,dx_j
\right).
\]
Pulling back, the second summand vanishes because $
f(\gamma(t))=0.
$
For the first summand, if \(i\neq i_0\), then \(f/x_i\) still contains the
factor \(x_{i_0}\), so
\[
\left(\frac{f}{x_i}\right)(\gamma(t))=0.
\]
If \(i=i_0\), then \(dx_{i_0}(t)=0\), because \(x_{i_0}(t)\equiv0\).
Therefore every term in \(\gamma^*\omega\) vanishes, and hence
$
\gamma^*\omega=0.
$
Thus every arc on \(\Dinv\) is tangential.

We have proved, after arbitrary extension of the ground field, that the
\(K[[t]]\)-points of the reduced tangential arc functor centred on
\(\Dinv\) coincide with the \(K[[t]]\)-points of \(J_\infty(\Dinv)\).  The
same argument applies after imposing the closed condition \(\gamma(0)\in Z\)
for any closed subset \(Z\subset\Dinv\).  Therefore
\[
J_\infty^{\tan}(W,\cG;Z)_{\red}
=
J_\infty(\Dinv;Z)_{\red}.
\]
 
\end{proof}


\begin{thebibliography}{99}
\bibitem[Carter]{Carter} \bibaut{P. J. Carter}, \emph{On the resolutions of non-dicritical foliations}, arXiv:2203.11094, 2024.
\bibitem[CarterThesis]{CarterThesis} \bibaut{P. J. Carter}, \emph{On the resolution of foliation singularities via jet schemes}, Ph.D. thesis, University of Liverpool, 2020.
\end{thebibliography}
\end{document}
